\section{Importance Discrimination Collapse under Compound Sparsity}
\label{sec:analysis}

The mechanisms discussed here are \emph{correlational diagnostics} that illuminate but do not fully explain the $15$--$17{\times}$ severity divergence between strategies. Both strategies exhibit near-identical Gini drops ($<10\%$ difference), yet severity diverges by $15$--$17{\times}$---an order-of-magnitude mismatch indicating that aggregate concentration collapse alone is insufficient. The critical differentiator is \emph{structural specificity} of importance-based pruning, which systematically removes high-attention-score heads (\S\ref{sec:discrimination-collapse-main}).

\subsection{Importance Discrimination Collapse}
\label{sec:discrimination-collapse-main}

Head pruning systematically reduces the discriminative power of cumulative importance scores, providing a correlational diagnostic for the phase-strategy interaction (causal attribution requires future intervention experiments). The Gini coefficient decreases by $0.030$--$0.081$ (DeepSeek, $6$--$17\%$ relative to dense baseline) and $0.010$--$0.020$ (Qwen, $2$--$4\%$) after pruning, with consistent negative deltas in all 16 conditions (Appendix~\ref{app:analysis}, Table~\ref{tab:discrimination-metrics-full}). Random eviction shows near-identical Gini drops (e.g., DeepSeek $s{=}0.25$: $-0.065$ vs.\ $-0.060$), indicating that collapse is a consequence of head pruning \emph{per se}. Yet importance-based severity is $15$--$17{\times}$ larger---the amplification arises because importance-based eviction \emph{depends on} the very scores that pruning degrades.

This pattern---\emph{importance discrimination collapse}---reflects eviction degenerating into noise-sensitive selection as score distributions approach uniformity. The effect is phase-dependent: per-phase Gini analysis shows that the low-sparsity phase exhibits the largest Gini delta at every sparsity level on both models, with monotonically increasing magnitude (DeepSeek: $-0.016$ at $s{=}0.15$, $-0.021$ at $s{=}0.20$, $-0.031$ at $s{=}0.25$, $-0.037$ at $s{=}0.30$; Qwen: $-0.017$, $-0.020$, $-0.029$, $-0.035$). Low-sparsity phases have broadly distributed attention where collapse has maximum impact; high-sparsity phases preserve concentration. \cite{hies-2025} report this collapse under single-axis pruning; we extend it to the compound setting.

To provide quantitative evidence for structural specificity, we analyzed the attention Gini coefficients of importance-pruned versus retained heads. Importance pruning (Frobenius norm criterion) systematically removes heads with lower Gini coefficients---i.e., heads with more dispersed attention patterns (Mann-Whitney $U$, $p < 10^{-28}$; rank-biserial $r = 0.62$--$0.76$ across all 8 conditions). Random pruning, by contrast, samples heads with Gini indistinguishable from the population mean (200-trial verification). The Frobenius norm and Gini share sensitivity to attention magnitude; the large effect size ($r > 0.6$) nonetheless indicates strong structural separation.

We note that this evidence is correlational; causal validation through controlled manipulation of importance-score entropy is an important direction for future work.

\subsection{Entropy Redistribution: A Model-Dependent Secondary Effect}
\label{sec:entropy-redistribution-main}

As a supplementary observation, head pruning also redistributes the attention entropy of surviving heads on DeepSeek ($\Delta\mathcal{H} \in [-0.047, -0.040]$; Appendix~\ref{app:analysis}, Table~\ref{tab:entropy-deltas}). However, entropy redistribution shows model-dependent behavior---the direction reverses on Qwen at $s{=}0.25$ ($\Delta\mathcal{H} = +0.045$)---replicating in only 2/4 Qwen conditions and thus only partially consistent. Importance discrimination collapse remains the more robust, model-independent correlate, though it partially accounts for rather than fully explains the severity divergence.
